\textbf{Staircase}

The equilibrium shape of bodies in zero-gravity is determined by the minimum of their surface energy. Thus, for example, the equilibrium shape of a water droplet turns out to be spherical: the sphere has the smallest surface area among bodies of the same volume. At low temperature, the equilibrium shape of crystals may have flat facets. The parts of the crystal surface that have a small angle $\phi$ with the facet are in fact staircases of rare steps on this facet. The height of such steps is equal to the period of the crystal lattice $h$.

\begin{center}\includegraphics[width=0.8\linewidth]{figures/EuPhO_2018_Q3_fig1.png}\end{center}

Equilibrium surface profile $y\left(x\right)$ of a certain crystal and the corresponding microscopic staircase are shown schematically in the figure, where n denotes the step number, counting from $x=0$. The profile shape at $x>0$ can be approximated as $y\left(x\right)=-\left(\frac{x}{\lambda}\right)^{\frac{3}{2}}h$ , where $\lambda=45\ \mu m$ and $h=0.3\ nm$.

a) Express the distance $d_n$ between two adjacent steps as a function of $n$ for $n\gg1$.

b) The interaction energy $E$ of two steps depend on the distance $d$ between them as
\[
E\left(d\right)=\mu d^{\nu}
\]
where $\mu$ is a constant. Assume that only adjacent steps interact. Find the numerical value for the exponent $\nu$.
